\section{引言与问题定义}
大模型推理中的参数访存与条件计算使传统“计算单元与存储单元分离”的组织方式面临带宽和容量压力。CIM 将权重驻留在计算邻近的存储阵列中，为降低数据移动提供了硬件抽象；已有加速器研究也表明，数据布局、复用关系和片上资源分配会共同决定实际吞吐\cite{sze2017efficient}。但在多 Sub-Cube 的 3D 空间中，权重块的物理位置又直接决定可并行度、切换次数和流水线等待。MoE 进一步放大了这一耦合：路由专家仅在被门控选中时执行，而共享专家或共享层具有稳定热点，因而仅按模型拓扑或体积进行静态放置难以稳定降低排队代价。

MoE 通过稀疏路由以较低的单次计算量激活较大的参数集合，Switch Transformer 和 GShard 已展示了这一范式在大规模模型训练中的可扩展性\cite{fedus2022switch,lepikhin2021gshard}。DeepSeek-V3 的公开技术报告也体现出细粒度路由专家与共享专家并存的结构特征\cite{deepseek2024v3}。这些工作说明，专家访问模式不是独立同分布的随机噪声，而是可由激活轨迹归纳的调度信息。

本项目聚焦赛题二而非完整的指令级编译器：输入为 JSON 或受限 ONNX 元数据与 MoE 激活轨迹，输出为权重空间映射和静态调度。目标是在固定的 3D CIM 资源配置内，首先最小化端到端推理周期，其次提升空间利用和大模型适配性；输入输出格式与约束以赛题原文为准\cite{contest2026}。本文的核心论点是：在不允许运行时迁移的前提下，将轨迹统计前置到静态映射，并将映射结果反馈给依赖感知调度，可以在当前资源模型和 mock trace 上减少资源等待与切换代价。本文的主要工作如下：
\begin{enumerate}
\item 建立模型解析、二维轴对齐 Weight-Cube 切分、轻量 JSON internal IR、映射和仿真的端到端可追溯链路；
\item 提出由共现、频率、burstiness 与转移统计驱动的 MoE 分组与副本放置策略，并结合热子图精确分配和局部搜索；
\item 实现考虑依赖屏障、Sub-Cube 互斥、切换、z 层访问、带宽和传输/计算重叠的静态调度模拟器，以及独立提交校验器；
\item 给出严格容量下的主方案、候选方案、复制消融和 \LargeExperts{} 专家 metadata-only 压力验证，并明确其数据与硬件模型边界。
\end{enumerate}

本文不把模拟周期解释为真实芯片实测频率，也不宣称完成完整 ONNX 执行语义、真实 DeepSeek-671B 权重部署或官方激活轨迹验证。所有结果均在实验章节中给出输入、配置、基线和验证状态。
